% Section 3.3, from lectures/L03/appendix/b_exercises.md.

\section{Exercises}
\label{c3:sec:exercises}\label{c3:app:b}
These exercises reinforce the concepts from \secref{c3:app:a}.

\begin{aside}
\textbf{How to check your work.} Write each exercise set as its own application under
\file{lectures/L03/exercises}, with the file names given there, and run \code{make test} from the
repository root. A set’s tests switch on as soon as its files exist, and a set you have not started
is reported as \code{SKIP}. Each application grows through its set, and the tests check where it
ends up. Programs run on \code{native_sim} in simulated time, so their output, timestamps included,
is compared with the expected output line by line.

The solutions are published in the companion repository, under
\file{lectures/L03/appendix/solutions}, a few days after the lecture. Try each exercise yourself
before you read them.
\end{aside}

Remember \code{source env.sh} in every new terminal. Build and run an exercise from the repository
root:

\begin{shell}
west build -p -b native_sim/native/64 lectures/L03/exercises/exercise1
./build/zephyr/zephyr.exe -stop_at=2 -no-rt
\end{shell}

A program that never stops by itself, even with \code{-stop_at}, is stuck in a busy loop: press
\code{Ctrl+C} \textbf{twice}.

% ---------------------------------------------------------------------
\exerciseset{Who runs when}
\label{c3:set:1}
Create an application named \code{exercise1} to work in for this exercise set, with a
\code{CMakeLists.txt} as in Chapter~\ref{c1:ch}, and \code{CONFIG_THREAD_NAME=y} in \code{prj.conf}.

Every line the program prints comes from this function, which prints the uptime, and the name and
priority of the thread that calls it. Copy it into \code{src/main.c}:

\begin{ccode}
static void trace(const char *what)
{
	k_tid_t self = k_current_get();

	printk("%4lld ms  %-7s (%2d)  %s\n", k_uptime_get(), k_thread_name_get(self),
	       k_thread_priority_get(self), what);
}
\end{ccode}

\exercise{Three threads and main}{Code}
\label{c3:ex:1.1}

\task{Tasks}
\begin{itemize}
  \item Write an entry function \code{periodic()} for a thread, which for now only calls
    \code{trace("started")}.
  \item Define three threads with \code{K_THREAD_DEFINE()}, all running \code{periodic()}, each with a stack
    of 1024 bytes, no options and no start delay:
\begin{narrowtable}{@{}ll@{}}
  \thead{Thread} & \thead{Priority} \\
  \midrule
  \code{control} & -1 \\
  \code{sensor} & 3 \\
  \code{logger} & 7 \\
\end{narrowtable}
  \item In \code{main()}, call \code{trace("started")} and return.
\end{itemize}

Before you run it, write down the order in which you expect the four lines. Then run it.

\task{Reflection}
\begin{itemize}
  \item In which order are the four lines printed? Why does \code{control} print before \code{main}, while
    \code{sensor} and \code{logger} print after it?
  \item Change \code{control}’s priority to 1 and run it again. What changes, and why?
\end{itemize}

\exercise{Periodic threads}{Code}
\label{c3:ex:1.2}

\task{Tasks}
Pass each thread two numbers through its first two arguments, with \code{INT_TO_POINTER()}, and read
them in \code{periodic()} with \code{POINTER_TO_INT()}:

\begin{narrowtable}{@{}lll@{}}
  \thead{Thread} & \thead{\code{p1}: period, in ms} & \thead{\code{p2}: work, in ms} \\
  \midrule
  \code{control} & 300 & 150 \\
  \code{sensor} & 200 & 0 \\
  \code{logger} & 100 & 0 \\
\end{narrowtable}

Extend \code{periodic()} so that, after \code{trace("started")}, it runs one round per period up to
600 ms:

\begin{ccode}
for (int round = 1; round * period_ms <= 600; round++) {
	k_msleep(period_ms);
	trace("awake");
	if (work_ms > 0) {
		k_busy_wait(work_ms * USEC_PER_MSEC);
		trace("work done");
	}
}
\end{ccode}

\code{k_busy_wait()} stands for work that keeps the CPU busy: unlike a loop, it lets simulated time
pass.

\task{Reflection}
\begin{itemize}
  \item At which uptimes does \code{logger} wake up? Why are they not multiples of 100?
  \item \code{logger}’s third round should start at 330 ms. When is it printed, and why?
\end{itemize}

\exercise{On the dot}{Code}
\label{c3:ex:1.3}

\task{Tasks}
Replace the relative sleep with an \textbf{absolute} one, so that each thread wakes at exactly every
multiple of its period, up to and including 600 ms. The loop becomes a loop over the next wake-up
time, \code{next}, from \code{period_ms} to 600 in steps of \code{period_ms}, with
\code{k_sleep(K_TIMEOUT_ABS_MS(next))} at the top of each round.

The output must be:

\begin{console}
   0 ms  control (-1)  started
   0 ms  main    ( 0)  started
   0 ms  sensor  ( 3)  started
   0 ms  logger  ( 7)  started
 100 ms  logger  ( 7)  awake
 200 ms  sensor  ( 3)  awake
 200 ms  logger  ( 7)  awake
 300 ms  control (-1)  awake
 450 ms  control (-1)  work done
 450 ms  sensor  ( 3)  awake
 450 ms  logger  ( 7)  awake
 460 ms  logger  ( 7)  awake
 500 ms  logger  ( 7)  awake
 600 ms  control (-1)  awake
 750 ms  control (-1)  work done
 750 ms  sensor  ( 3)  awake
 750 ms  logger  ( 7)  awake
\end{console}

\textbf{Hint:} \code{next} must be an \code{int64_t}, like the uptime.

\task{Reflection}
\begin{itemize}
  \item Which of \code{logger}’s rounds is printed at 450 ms, which at 460 ms, and why is the second one 10
    ms later rather than at the same moment?
  \item Which round of \code{sensor} is printed at 450 ms?
  \item \code{control} is cooperative. Make it preemptible, with priority 1. Does anything change between
    300 and 450 ms? Why, or why not?
  \item Keep \code{control} at priority 1 and give \code{sensor} priority 0. What happens at 400 ms now? Why
    is \code{control}’s work still done at 450 ms, although it was interrupted?
  \item Replace \code{k_busy_wait()} with a \code{for} loop that counts to a few hundred million. What
    happens to the trace on \code{native_sim}, and why would it be different on a microcontroller?
\end{itemize}

% ---------------------------------------------------------------------
\exerciseset{A timer and the system work queue}
\label{c3:set:2}
Create an application named \code{exercise2} to work in for this exercise set, with
\code{CONFIG_THREAD_NAME=y} in \code{prj.conf}.

Use this function to name the context that code runs in:

\begin{ccode}
static const char *context(void)
{
	return k_is_in_isr() ? "ISR" : k_thread_name_get(k_current_get());
}
\end{ccode}

\exercise{From ISR to thread}{Code}
\label{c3:ex:2.1}

\task{Tasks}
\begin{itemize}
  \item Define a work item \code{sample_work} with \code{K_WORK_DEFINE()}, whose handler
    \code{sample_handler()} counts the samples it has taken.
  \item Define a timer \code{sample_timer} with \code{K_TIMER_DEFINE()}, whose expiry function
    \code{timer_expired()} stores \code{context()} in a variable \code{expiry_context} and submits
    \code{sample_work} to the system work queue.
  \item In \code{main()}, start the timer with a delay of 250 ms and a period of 250 ms, then print
    \code{timer started in <context>}, prefixed with the uptime as in the output below, and return.
  \item In \code{sample_handler()}, print a line for each sample with its number, \code{expiry_context} and
    \code{context()}, and stop the timer with \code{k_timer_stop()} after the fourth sample.
\end{itemize}

Every line starts with the uptime, formatted with \code{"\%4lld ms  "}. Exercise~\ref{c3:ex:2.2}
gives the lines their final form.

\exercise{What a submission does, and a stop function}{Code}
\label{c3:ex:2.2}

\task{Tasks}
\begin{itemize}
  \item In \code{timer_expired()}, submit \code{sample_work} \textbf{twice}, and store the two return values
    of \code{k_work_submit()}.
  \item Give the timer a stop function, \code{timer_stopped()}, which prints \code{timer stopped, stop
    function in <context>}.
  \item Print the return values and both contexts in \code{sample_handler()}.
\end{itemize}

The output must be:

\begin{console}
   0 ms  timer started in main
 250 ms  sample 1: expiry in ISR, submitted 1 then 0, handler in sysworkq
 500 ms  sample 2: expiry in ISR, submitted 1 then 0, handler in sysworkq
 750 ms  sample 3: expiry in ISR, submitted 1 then 0, handler in sysworkq
1000 ms  sample 4: expiry in ISR, submitted 1 then 0, handler in sysworkq
1000 ms  timer stopped, stop function in sysworkq
\end{console}

\task{Reflection}
\begin{itemize}
  \item The work was submitted twice for every expiry. How many times did the handler run, and what do the
    return values 1 and 0 mean?
  \item Why does the stop function run in \code{sysworkq}, and not in an ISR like the expiry function?
  \item The samples come at 250, 500, 750 and 1000 ms, while \code{k_msleep(250)} in a loop would wake at
    260, 520, and so on. Why do the timer’s expiries not drift?
  \item Add \code{k_msleep(10)} to \code{timer_expired()} and run it. Then build again with \code{--
    -DCONFIG_ASSERT=y} and run it. What happens each time, and what does that tell you about assertions?
  \item Remove the \code{k_msleep(10)}, and instead add \code{k_msleep(300)} to \code{sample_handler()},
    after the print and before the check that stops the timer, so that handling a sample takes longer
    than the period. When are the samples taken now? What does \code{k_work_submit()} return, and why is
    there a fifth sample, after the timer was stopped?
\end{itemize}

% ---------------------------------------------------------------------
\exerciseset{Debouncing}
\label{c3:set:3}
Create an application named \code{exercise3} to work in for this exercise set.

A mechanical button does not change level once when pressed: its contacts bounce, and the pin
toggles several times over a few milliseconds. \textbf{Debouncing} means reporting a new level only
when the pin has stayed at it for a while. This set uses Chapter~\ref{c2:ch}’s emulated GPIO
controller, and a thread that plays a finger with bouncing contacts.

Use these files as they are. \code{prj.conf}:

\begin{kconfigcode}
# The GPIO API, for the button.
CONFIG_GPIO=y

# A tick every millisecond, so that a bounce a few milliseconds long can be timed.
CONFIG_SYS_CLOCK_TICKS_PER_SEC=1000
\end{kconfigcode}

\code{app.overlay}, a button on pin 1 of the emulated GPIO controller:

\begin{dtscode}
#include <zephyr/dt-bindings/gpio/gpio.h>

/ {
	aliases {
		sw0 = &button0;
	};

	buttons {
		compatible = "gpio-keys";

		button0: button_0 {
			gpios = <&gpio0 1 (GPIO_PULL_UP | GPIO_ACTIVE_LOW)>;
			label = "Button";
		};
	};
};
\end{dtscode}

And in \code{src/main.c}, the finger: a press that bounces, a release that bounces, and a glitch too
short to be a press. \code{edges} is the count of edges your callback has seen, below.

\begin{ccode}
/* Set the button's physical level: 0 is pressed, since the button is active low. */
static void level(int value, int hold_ms)
{
	gpio_emul_input_set(button.port, button.pin, value);
	k_msleep(hold_ms);
}

static void finger(void *p1, void *p2, void *p3)
{
	k_sleep(K_TIMEOUT_ABS_MS(100));
	level(0, 2);
	level(1, 1);
	level(0, 3);
	level(1, 1);
	level(0, 300);

	level(1, 4);
	level(0, 2);
	level(1, 300);

	level(0, 5);
	level(1, 300);

	printk("%4lld ms  finger done, %ld edges\n", k_uptime_get(), (long)atomic_get(&edges));
}

K_THREAD_DEFINE(finger_thread, 1024, finger, NULL, NULL, NULL, 7, 0, 0);
\end{ccode}

\exercise{Every edge}{Code}
\label{c3:ex:3.1}

\task{Tasks}
\begin{itemize}
  \item Get the button as a \code{struct gpio_dt_spec} named \code{button}, from the \code{sw0} alias.
  \item In \code{main()}, check that the button is ready, configure it as an input, configure an interrupt
    on \textbf{both} edges, and add a callback.
  \item In the callback, count every edge in an \code{atomic_t} named \code{edges}.
\end{itemize}

Run it: the finger reports 10 edges at 1028 ms.

\exercise{Only real presses}{Code}
\label{c3:ex:3.2}

\task{Tasks}
\begin{itemize}
  \item Define a delayable work item, \code{debounce_work}, with a handler that reads the button’s logical
    level with \code{gpio_pin_get_dt()}, and prints a line only when that level differs from the last
    one it reported (the button starts released):
\begin{console}
<uptime> ms  pressed, after <edges> edges
<uptime> ms  released, after <edges> edges
\end{console}
  \item In the callback, after counting the edge, (re)schedule \code{debounce_work} to run 20 ms after the
    \textbf{last} edge.
\end{itemize}

The output must be:

\begin{console}
 132 ms  pressed, after 5 edges
 441 ms  released, after 8 edges
1028 ms  finger done, 10 edges
\end{console}

\textbf{Hint:} Of \code{k_work_schedule()} and \code{k_work_reschedule()}, only one moves an already
scheduled deadline.

\task{Reflection}
\begin{itemize}
  \item The last edge of the press is at 111 ms. Why is the press reported at 132 ms, and not at 131?
  \item Why is the glitch at 721 ms, two edges, never reported?
  \item Try the other one of \code{k_work_schedule()} and \code{k_work_reschedule()}. When are the press and
    the release reported now? The output still has the same events: describe a finger for which it would
    not.
  \item On \code{native_sim}, the callback runs in the finger’s thread. In which context would it run on a
    microcontroller, and why is it still correct to call your scheduling function from it?
  \item Only the work handler reads and writes the variable holding the last reported level, so it needs no
    lock. Why can the handler not run twice at the same time?
  \item Set \code{CONFIG_SYS_CLOCK_TICKS_PER_SEC=100} and run it. When are the press and the release
    reported, and why?
\end{itemize}

% ---------------------------------------------------------------------
\exerciseset{Stop spinning}
\label{c3:set:4}
Create an application named \code{exercise4} to work in for this exercise set, with a
\code{prj.conf} holding only a comment.

This program reads three channels of a converter, whose conversion takes 40 ms, simulated with a
timer. Copy it into \code{src/main.c}:

\begin{ccode}
#include <stdbool.h>

#include <zephyr/kernel.h>

/* How long the simulated conversion takes. */
#define CONVERSION_MS 40

/* Set by the timer's expiry function when the conversion is done. */
static volatile bool conversion_done;

static void conversion_expired(struct k_timer *timer)
{
	conversion_done = true;
}

K_TIMER_DEFINE(conversion_timer, conversion_expired, NULL);

/* Start a conversion on a channel, wait for it to finish, and return the result. */
static int convert(int channel)
{
	conversion_done = false;
	k_timer_start(&conversion_timer, K_MSEC(CONVERSION_MS), K_NO_WAIT);

	while (!conversion_done) {
	}

	return 1000 + channel * 111;
}

int main(void)
{
	for (int channel = 0; channel < 3; channel++) {
		const int value = convert(channel);

		printk("%3lld ms  channel %d: %d\n", k_uptime_get(), channel, value);
	}
	printk("%3lld ms  all channels converted\n", k_uptime_get());

	return 0;
}
\end{ccode}

\exercise{The stall}{Code}
\label{c3:ex:4.1}

\task{Tasks}
Build it, and run it with \code{./build/zephyr/zephyr.exe -stop_at=1 -no-rt}. Then run \code{make
test}.

\task{Reflection}
\begin{itemize}
  \item Why does the program print nothing, not even channel 0, and why does \code{-stop_at=1} not stop it?
  \item What happens when you press \code{Ctrl+C} once? And twice?
  \item What does \code{make test} report for the set, and after how long?
  \item Would this program work on a microcontroller? Would it be a good program there?
  \item Would \code{k_yield()} inside the loop fix it on \code{native_sim}? Try it, and explain.
\end{itemize}

\exercise{A semaphore}{Code}
\label{c3:ex:4.2}

\task{Tasks}
Replace the flag with a semaphore: the expiry function gives it, and \code{convert()} takes it with
\code{K_FOREVER} instead of spinning.

The output must be:

\begin{console}
 40 ms  channel 0: 1000
 80 ms  channel 1: 1111
120 ms  channel 2: 1222
120 ms  all channels converted
\end{console}

\exercise{No semaphore at all}{Code}
\label{c3:ex:4.3}

\task{Tasks}
Remove the semaphore and the expiry function as well, define the timer with
\code{K_TIMER_DEFINE(conversion_timer, NULL, NULL)}, and wait for it in \code{convert()} with
\code{k_timer_status_sync()}. The output must be unchanged.

\task{Reflection}
\begin{itemize}
  \item Why do the conversions finish at exactly 40, 80 and 120 ms?
  \item Suppose the timer expired before \code{convert()} started to wait. Would the semaphore version still
    work? Would the \code{k_timer_status_sync()} version?
  \item Zephyr’s documentation of \code{native_sim} suggests calling \code{k_cpu_idle()} inside a busy loop
    that cannot be removed. Try it in the original program. Does it work, and why is 4.3 still the
    better fix?
\end{itemize}

% ---------------------------------------------------------------------
\exerciseset{Stacks without a stack analyzer}
\label{c3:set:5}
No code in this set: answer the questions by building and looking.

\exercise[fillaccent2]{Frames}{Reflection}
\label{c3:ex:5.1}

\task{Tasks}
Build your \code{exercise2} with static stack usage analysis:

\begin{shell}
west build -p -b native_sim/native/64 lectures/L03/exercises/exercise2 -- -DCONFIG_STACK_USAGE=y
\end{shell}

\noindent Answer, using \code{build/CMakeFiles/app.dir/src/main.c.su} and the other \code{.su} files
in \code{build/}:

\begin{itemize}
  \item How large is \code{sample_handler()}’s stack frame?
  \item On which thread’s stack does \code{sample_handler()} run, and how large is that stack? Find the
    option that sets it in \code{build/zephyr/.config}.
  \item Is \code{sample_handler()}’s frame all it needs? Which functions does it call, and how large are
    their frames?
  \item On a microcontroller, on which stack would \code{timer_expired()} run, and which option sets its
    size?
\end{itemize}

\exercise[fillaccent2]{What does not work here}{Reflection}
\label{c3:ex:5.2}

\task{Tasks}
Build \code{exercise2} again with \code{-- -DCONFIG_THREAD_ANALYZER=y}.

\task{Reflection}
\begin{itemize}
  \item What happens, and which dependency is not met?
  \item Why can a thread analyzer that measures how much of each thread’s stack was used not work on
    \code{native_sim}?
\end{itemize}
